% =============================================================================
% SECTION 5: FINDINGS
% =============================================================================
\section{Findings}
\label{sec:findings}

This section synthesizes the evidence from 160 primary studies to answer our five research
questions. For each RQ, we present the aggregate finding, supporting evidence from key studies,
and a critical analysis of the evidence quality.

\subsection{RQ1: What Paradigms Have Been Proposed?}
\label{subsec:rq1}

\paragraph{Finding.}
\textbf{Seven distinct repository representation paradigms (\parcmd{1}--\parcmd{7}) have been
proposed in the literature, with PAR-2 (RAG) being the most prevalent (48/160 studies, 30\%)
and PAR-7 (Agentic) showing the fastest growth rate (from 2\% in 2022 to 18\% in 2025--2026).}

The seven-paradigm taxonomy presented in Section~\ref{sec:taxonomy} accounts for all 160
primary studies without remainder. No study fell outside the taxonomy, confirming that our
inductive synthesis has achieved saturation. The distribution across paradigms reveals a highly
skewed landscape: the top three paradigms (\parcmd{2}, \parcmd{3}, \parcmd{7}) account for 62\%
of studies, while \parcmd{6} (Memory) is the least-represented pure paradigm with only 12
studies.

The historical development of paradigms follows a recognizable pattern of increasing
sophistication. The first generation (2020--2022) focused on representation quality through
dense embeddings~\cite{feng2020codebert, guo2021graphcodebert, guo2022unixcoder}. The second
generation (2022--2023) focused on context augmentation through retrieval~\cite{zhang2023repocoder,
lu2022reacc, shrivastava2023repolevel}. The third generation (2024--2026) focuses on
structural richness through graphs~\cite{liu2024graphcoder, ouyang2025repograph, tao2025cgm}
and adaptive intelligence through agents~\cite{yang2024sweagent, zhang2024autocoderover,
ma2025lingma}. This progression reflects not only technical maturation but also the expanding
scope of SE tasks targeted: as researchers moved from code completion (T1) to issue
resolution (T2, T6) and feature implementation (T3), the representational demands increased
accordingly.

\paragraph{Evidence Quality.}
The evidence for this finding is based on the full set of 160 studies and is supported by the
inter-rater reliability analysis reported in Section~\ref{subsec:extraction}. The main
uncertainty is whether the taxonomy will remain stable as the field evolves: it is possible
that post-2026 developments (e.g., natively multi-file context LLMs with context windows of
10M+ tokens) will render some paradigms obsolete or merge currently distinct paradigms.

\subsection{RQ2: How Do Paradigms Map Onto SE Tasks?}
\label{subsec:rq2}

\paragraph{Finding.}
\textbf{The task--paradigm relationship is not random: structured tasks with high locality
(T1, T7) are dominated by simpler paradigms (\parcmd{1}, \parcmd{2}), while complex tasks
requiring global reasoning (T2, T6) attract diverse paradigm mixtures including \parcmd{4}
and \parcmd{7}.}

As detailed in Section~\ref{sec:semapping}, code completion (T1) is overwhelmingly addressed
by \parcmd{2} (RAG), while bug localization (T2) and program repair (T6) show the most diverse
paradigm distributions. This pattern reflects an underlying task property: code completion is
locally constrained (the ground truth completion depends on a bounded neighborhood of the
cursor), while bug localization and repair are globally unconstrained (the root cause may be
in any part of the repository). Globally unconstrained tasks require exploration strategies
that can dynamically focus attention, which is exactly the strength of \parcmd{7} (Agentic)
and \parcmd{3} (graph-guided traversal).

Test generation (T5) presents an interesting case. Despite being a globally unconstrained task
(generating correct tests requires understanding all dependencies of the unit under test),
it is primarily addressed by \parcmd{4} (Dynamic) and \parcmd{7} (Agentic), with notably few
applications of \parcmd{3} (Graphs). We hypothesize that the test generation community has
not yet fully adopted graph-based dependency analysis, which could provide a principled
basis for generating dependency-aware test stubs and mocks. This represents a significant
gap identified in Section~\ref{sec:semapping}.

Code translation (T4) is an emerging task with a distinctive profile: it requires both global
repository structure understanding (to map inter-module dependencies from source to target
language) and local idiom translation (to convert individual code constructs). No paradigm
has established dominance for T4, suggesting that this is an active frontier where hybrid
approaches~\cite{ibrahimzada2025alphatrans, feng2026dependency} are likely to yield the
greatest advances.

\subsection{RQ3: What Are the Empirical Performance Characteristics?}
\label{subsec:rq3}

\paragraph{Finding.}
\textbf{Graph-augmented approaches consistently outperform pure RAG for structured tasks (bug
localization, impact analysis) by 5--15 percentage points; agentic frameworks dominate SWE-bench
issue resolution with top systems achieving 40--50\% pass@1; dense embedding quality has
plateaued with marginal improvements from 2023 onward.}

\subsubsection{Code Completion Performance}

Table~\ref{tab:completion-perf} summarizes representative performance results on RepoBench
and CrossCodeEval for code completion across paradigms.

\begin{table*}[ht]
\caption{Representative code completion performance on RepoBench (EM\%) and CrossCodeEval (EM\%).}
\label{tab:completion-perf}
\small
\begin{tabular}{lllrr}
\toprule
\textbf{System} & \textbf{Primary PAR} & \textbf{Base LLM} & \textbf{RepoBench EM} & \textbf{XCodeEval EM} \\
\midrule
RepoCoder~\cite{zhang2023repocoder}           & PAR-2 & CodeT5        & 23.4 & -- \\
CrossCodeEval~\cite{ding2023crosscodeeval}    & PAR-1 & CodeGen       & -- & 18.2 \\
RepoHyper~\cite{phan2025repohyper}            & PAR-3 & StarCoder     & 31.7 & 24.1 \\
GraphCoder~\cite{liu2024graphcoder}           & PAR-3 & DeepSeek-Coder & 34.2 & 25.8 \\
HyRACC~\cite{li2025hyracc}                   & PAR-2 & GPT-4         & 36.1 & 27.3 \\
AlignCoder~\cite{jiang2025aligncoder}         & PAR-2 & Code Llama    & 33.8 & 25.1 \\
CodeRAG~\cite{zhang2025coderag}               & PAR-2 & GPT-4         & 38.4 & 29.2 \\
GRACE~\cite{wang2025grace}                    & PAR-3 & DeepSeek-Coder & 37.1 & 28.6 \\
\bottomrule
\end{tabular}
\end{table*}

The data in Table~\ref{tab:completion-perf} illustrates several important trends. First, graph
augmentation (\parcmd{3}) consistently outperforms pure RAG (\parcmd{2}) when the base LLM is
equivalent, suggesting that structural context genuinely improves completion quality beyond
what semantic similarity retrieval provides. Second, the choice of base LLM has a larger
impact on absolute performance than the choice of paradigm: switching from CodeT5 to GPT-4
with the same RAG pipeline improves EM by approximately 15 percentage points, while switching
from \parcmd{2} to \parcmd{3} with the same LLM improves EM by approximately 3--5 percentage
points. Third, there is notable variance across benchmarks: systems optimized for RepoBench
do not always generalize to CrossCodeEval, reflecting differences in the linguistic distribution
and difficulty distribution of the two benchmarks.

\subsubsection{Issue Resolution Performance}

SWE-bench~\cite{jimenez2024swebench} has become the de facto benchmark for repository-level
issue resolution, making it possible to compare systems across studies. Table~\ref{tab:swebench-perf}
summarizes representative results on the SWE-bench Verified subset.

\begin{table}[h]
\caption{Representative performance on SWE-bench Verified (pass@1 \%).}
\label{tab:swebench-perf}
\small
\begin{tabular}{llr}
\toprule
\textbf{System} & \textbf{Primary PAR} & \textbf{Pass@1 (\%)} \\
\midrule
SWE-agent~\cite{yang2024sweagent}           & PAR-7 & 18.1 \\
AutoCodeRover~\cite{zhang2024autocoderover} & PAR-7 & 22.7 \\
Agentless~\cite{xia2025demystifying}        & PAR-2 & 27.3 \\
LingmaAgent~\cite{ma2025lingma}             & PAR-7 & 33.6 \\
SWE-RL~\cite{wei2025swerl}                  & PAR-4 & 41.0 \\
SWE-Gym~\cite{pan2025training}              & PAR-7 & 38.2 \\
\bottomrule
\end{tabular}
\end{table}

The rapid progress on SWE-bench since its introduction in 2024 is striking: pass@1 has risen
from under 5\% for the earliest evaluated systems to over 40\% for the best current systems.
PAR-7 (Agentic) and PAR-4 (Dynamic/execution) dominate this benchmark, consistent with the
task's requirement for flexible, multi-step reasoning. Notably, Agentless~\cite{xia2025demystifying}
achieves competitive results with a non-agentic PAR-2 approach, suggesting that aggressive
retrieval can partially compensate for the lack of iterative exploration.

\subsubsection{Code Search Performance}

For code search (T7), the primary metric is Mean Reciprocal Rank (MRR) on the CodeSearchNet
benchmark. Dense embedding systems (\parcmd{1}) have improved from MRR$\approx$0.71
(CodeBERT~\cite{feng2020codebert}) to MRR$\approx$0.84 (GraphCodeBERT~\cite{guo2021graphcodebert},
UniXcoder~\cite{guo2022unixcoder}) with incremental improvements through 2023--2025.
The plateau observed in recent work reflects the approaching ceiling of what can be achieved
with function-level embeddings on CodeSearchNet; researchers have shifted attention to more
challenging code search benchmarks that involve cross-file and cross-repository queries.

\subsection{RQ4: What Are the Key Trade-offs?}
\label{subsec:rq4}

\paragraph{Finding.}
\textbf{No paradigm dominates across all five evaluation dimensions; each paradigm occupies a
distinct region of the trade-off space, making paradigm selection an engineering decision that
depends on deployment constraints.}

Table~\ref{tab:tradeoffs} presents the five-dimensional trade-off analysis, evaluating each
PAR on a three-point scale (Low/Medium/High) for each dimension.

\begin{table*}[ht]
\caption{Five-dimensional trade-off analysis of repository representation paradigms.}
\label{tab:tradeoffs}
\small
\begin{tabular}{lp{2.2cm}p{2cm}p{2cm}p{2cm}p{2.5cm}}
\toprule
\textbf{PAR} & \textbf{Context Coverage} & \textbf{Retrieval Latency} & \textbf{Training Cost}
             & \textbf{Scalability} & \textbf{Benchmark Avail.} \\
\midrule
\parcmd{1} Dense Emb.  & Low--Med & Low ({\textless}100ms) & Med-High & High & High (CodeSN, etc.) \\
\parcmd{2} RAG         & Med & Low--Med & Low & High & High (RepoBench) \\
\parcmd{3} Graph       & Med--High & Med (1--10s build) & Med & Med & Med \\
\parcmd{4} Dynamic     & High (runtime) & High ({\textgreater}10s exec) & Med & Low & Low \\
\parcmd{5} Summ.       & High (compressed) & Low & Low & High & Low \\
\parcmd{6} Memory      & High (episodic) & Med & Low & Med & Low \\
\parcmd{7} Agentic     & High (adaptive) & High ({\textgreater}30s) & Med & Low & High (SWE-bench) \\
\bottomrule
\end{tabular}
\end{table*}

\paragraph{Context Coverage.}
Context coverage measures the fraction of repository-wide relevant information that the
paradigm can plausibly surface within a fixed context window. Dynamic (\parcmd{4}), Agentic
(\parcmd{7}), and Memory-augmented (\parcmd{6}) paradigms achieve the highest coverage because
they are not limited to a single retrieval step: they can iteratively gather information as
needed. Dense embeddings (\parcmd{1}) have the lowest coverage because a single nearest-neighbor
search has no mechanism to retrieve transitively relevant fragments.

\paragraph{Retrieval Latency.}
Retrieval latency is critical for interactive applications. \parcmd{1} (dense embedding search)
is the fastest, typically completing in under 100 milliseconds. \parcmd{7} (Agentic) is the
slowest, often requiring dozens of LLM calls totaling 30 seconds to several minutes. This
latency difference drives a hard partition in the application space: interactive IDE tools
must use \parcmd{1} or \parcmd{2}, while offline batch tasks (nightly CI, pre-commit checks)
can afford \parcmd{7}. The literature largely confirms this partition: all production IDE
integration studies use \parcmd{1} or \parcmd{2}~\cite{yang2025deepdive, deng2025enhancing,
li2024enhancing}.

\paragraph{Training Cost.}
Training cost encompasses both the computational resources required to train the representation
model and the maintenance cost of keeping it updated as the repository evolves. \parcmd{2}
(RAG) has the lowest training cost for a given LLM because it requires no task-specific
training at all---the retriever can be deployed immediately on a new repository by indexing
its content. \parcmd{1} (Dense Embeddings) incurs a moderate training cost for the initial
embedding model fine-tuning, but this cost is amortized across all downstream uses. \parcmd{4}
(Dynamic) incurs moderate training cost for RL-based systems such as IRCoCo and SWE-RL.

\paragraph{Scalability.}
Scalability measures how gracefully a paradigm handles repositories of increasing size (lines
of code) and complexity (number of interacting modules). \parcmd{1} and \parcmd{2} scale
well: the index can be partitioned across multiple machines, and retrieval time grows
logarithmically with index size using approximate nearest-neighbor algorithms. \parcmd{4}
(Dynamic) scales poorly because test execution time grows linearly or superlinearly with
repository size. \parcmd{7} (Agentic) scales poorly for the same reason: as repository size
grows, the number of exploration steps required grows, and the agent's probability of choosing
the wrong exploration path increases.

\paragraph{Benchmark Availability.}
Benchmark availability reflects the degree to which a paradigm's progress can be tracked with
standardized, community-accepted benchmarks. \parcmd{1} benefits from the mature CodeSearchNet
benchmark and the extensive HumanEval-inspired family. \parcmd{2} benefits from RepoBench
and CrossCodeEval. \parcmd{7} benefits from SWE-bench. In contrast, \parcmd{4}, \parcmd{5},
and \parcmd{6} lack well-established benchmarks: most studies in these paradigms evaluate on
bespoke datasets that are difficult to compare across studies.

\subsection{RQ5: Open Challenges and Future Directions}
\label{subsec:rq5}

\paragraph{Finding.}
\textbf{We identify seven critical open challenges in repository representation for LLM-based
software engineering, spanning benchmark standardization, hybrid paradigm design, energy
efficiency, temporal consistency, multilingual support, adversarial robustness, and
human-AI collaboration.}

\subsubsection{Challenge 1: Longitudinal and Evolutionary Benchmarks}

All major existing benchmarks (RepoBench, CrossCodeEval, SWE-bench) use static snapshots of
repositories. Real software development is an evolutionary process: the repository changes
continuously, and the LLM-based tool must maintain a consistent representation across these
changes. No benchmark currently evaluates the ability of a system to update its representation
efficiently as the repository evolves. Zhang et al.~\cite{zhang2025swebenchlive} make a step
toward this with SWE-bench Live, which uses continuously crawled GitHub issues, but the
representation update problem remains under-studied. HumanEvo~\cite{zheng2025humanevo} is
another nascent effort in this direction, evaluating code generation on evolving repositories.

\subsubsection{Challenge 2: Multi-Paradigm Hybridization}

As noted in Section~\ref{subsec:hybrid}, hybridization is increasingly common in 2025--2026.
However, there is no principled framework for determining when and how to combine paradigms.
Current hybrid systems are designed heuristically, and the interaction effects between
paradigms are poorly understood. A rigorous framework for paradigm composition---analogous
to ensemble methods in machine learning---would enable more systematic exploration of the
hybrid design space.

\subsubsection{Challenge 3: Energy-Aware Evaluation}

No study in our corpus reports energy consumption or carbon footprint as an evaluation metric.
As LLM-based software engineering tools are deployed at scale---GitHub Copilot alone is used
by millions of developers daily---their energy footprint becomes a significant environmental
concern. Future work should include standardized measurement of FLOPs, kWh, and
$\text{CO}_2$e as primary metrics alongside performance.

\subsubsection{Challenge 4: Multilingual and Multi-Stack Repositories}

The majority of studies (approximately 68\%) evaluate on Python-only repositories. Modern
software projects routinely mix multiple programming languages: a backend in Python, a frontend
in TypeScript, infrastructure in Terraform, tests in a mix of languages, and documentation
in Markdown and RST. Repository representation for polyglot projects is substantially harder
than for monolingual ones, and no existing paradigm has been systematically evaluated in
this setting. OmniGIRL~\cite{guo2025omnigirl} and Multi-SWE-bench~\cite{zan2025multiswebench}
begin to address this gap.

\subsubsection{Challenge 5: Adversarial Robustness}

LLM-based agents that represent repository content are potentially vulnerable to adversarial
manipulation: an attacker who can insert malicious code comments, misleading function names,
or carefully crafted API stubs into a repository could potentially hijack the agent's behavior.
This is a novel attack surface that has not yet been systematically studied. Given the
increasing deployment of LLM-based code agents in security-critical settings (automated
code review, CI/CD pipelines), adversarial robustness of repository representation deserves
urgent attention.

\subsubsection{Challenge 6: Human-AI Collaborative Representation}

Current paradigms treat repository representation as a fully automated process. However,
developers have rich mental models of their repositories that could substantially improve
representation quality if properly elicited. Interactive and collaborative paradigms---where
the LLM agent queries the developer for clarification, and the developer's responses are
incorporated into the repository representation---are largely unexplored. The RECODE-H
benchmark~\cite{miao2026recodeh} is an early step in this direction, simulating interactive
human feedback in research code development.

\subsubsection{Challenge 7: Formal Verification Integration}

Program repair and code generation tasks have correctness requirements that cannot be fully
captured by test-based evaluation. A patch that passes all tests may still introduce subtle
semantic bugs. Integrating formal verification tools (model checkers, theorem provers,
SMT solvers) into the repository representation pipeline could provide stronger correctness
guarantees, but the interface between LLM-generated code and formal verification tools
remains largely unexplored.
